\begin{lemma}[Vanishing of the higher derivatives at a critical point]
Let $1<p<2$, suppose $p$ satisfies Lewis rigidity, and let $(\kappa,C_0)$ be a first-order
estimate datum for $p$.  Let $\Omega\subseteq\mathrm{Plane}$, let
$u\colon\mathrm{Plane}\to\mathbb{R}$, let $x_0\in\Omega$ and let $m\in\mathbb{N}$ with
$m\ge2$.  Assume
\begin{enumerate}
\item $u$ is of class $C^{m}$ on $\Omega$;
\item $u$ is $p$-harmonic in $\Omega$ with weak gradient equal to its own classical gradient
$Du$;
\item $Du(x_0)=0$.
\end{enumerate}
Then
\[
  D^{j}u(x_0)[w,\ldots,w]=0\qquad\text{for every }j\text{ with }2\le j\le m
  \text{ and every }w\in\mathrm{Plane},
\]
where $D^{j}u(x_0)$ is the $j$-th Fr\'echet derivative of $u$ at $x_0$ and the brackets denote
its evaluation on $j$ copies of $w$.
\end{lemma}
\begin{proof}
Hypothesis (2) says that $u$ is $p$-harmonic in $\Omega$ with weak gradient $Du$; this is the
faithful rendering of the paper's hypothesis, since for a $C^{m}$ function the paper's $Du$ is
the classical gradient.  By condition (1) of \noderef{IsPHarm}, $\Omega$ is open, so, as
$x_0\in\Omega$, we may fix once and for all
\[
  \rho>0\qquad\text{with}\qquad B_\rho(x_0)\subseteq\Omega .
\]
The ball $B_\rho(x_0)$ is open, so part (a) of \noderef{PHarmRestrict} applied to hypothesis (2)
gives
\begin{equation}\label{eq:rig-ball}
  u\ \text{is}\ p\text{-harmonic in}\ B_\rho(x_0)\ \text{with weak gradient}\ Du ,
\end{equation}
and hypothesis (1) restricts to: $u$ is of class $C^{m}$ on $B_\rho(x_0)$.

For $1\le q\le m+1$ let $T(q)$ denote the statement
\[
  T(q):\qquad D^{j}u(x_0)[w,\ldots,w]=0\quad\text{for every }j\text{ with }1\le j<q
  \text{ and every }w\in\mathrm{Plane}.
\]
The conclusion of the lemma is contained in $T(m+1)$, so it suffices to prove $T(q)$ for every
$q$ with $1\le q\le m+1$.  We do this by induction on $q$.

\medskip\noindent\emph{Base cases.}
$T(1)$ is vacuous, there being no $j$ with $1\le j<1$.  For $T(2)$ the only index is $j=1$, and
by clause (6) of the background proposition on the recursion for the iterated Fr\'echet
derivative, $D^{1}u(x_0)[w]=\langle Du(x_0),w\rangle$ for every $w$, which vanishes by
hypothesis (3).  (The function $u$ is differentiable at $x_0$ because it is of class $C^{m}$ on
the open set $\Omega$ with $m\ge2\ge1$, by clause (2) of the background proposition on locality
of $C^{1}$ regularity.)

\medskip\noindent\emph{Inductive step.}
Let $2\le q\le m$ and assume $T(q)$.  Since $T(q+1)$ adds to $T(q)$ exactly the index $j=q$, it
suffices to prove
\begin{equation}\label{eq:rig-goal}
  D^{q}u(x_0)[w,\ldots,w]=0\qquad\text{for every }w\in\mathrm{Plane}.
\end{equation}
Define the continuous $q$-linear form and the degree-$q$ homogeneous polynomial
\[
  Q:=\frac{1}{q!}\,D^{q}u(x_0),\qquad P(x):=Q[x,\ldots,x]=\frac{1}{q!}\,D^{q}u(x_0)[x,\ldots,x].
\]
That $Q$ is a continuous $q$-linear form, and that $P$ is continuous and satisfies
$P(tx)=t^{q}P(x)$, are clauses (7) and (5) of the recursion background proposition; the
homogeneity uses multilinearity alone and no symmetry of $Q$.

\emph{Step 1: the blow-ups.}  For $r>0$ put
\[
  u_r(x):=\frac{u(x_0+r\,x)-u(x_0)}{r^{q}} .
\]
Apply \noderef{PHarmNormalize} to \eqref{eq:rig-ball}, with domain $B_\rho(x_0)$, functions $u$
and $Du$, centre $x_0$, dilation factor $r$, normalizing factor $M:=r^{q}>0$ and additive
constant $c:=u(x_0)$.  Its hypotheses hold ($1<p<2$, $r>0$, $r^{q}>0$), and its conclusion is
that $u_r$ is $p$-harmonic in $\{x:x_0+r\,x\in B_\rho(x_0)\}$ with weak gradient
$x\mapsto(r/r^{q})\,Du(x_0+r\,x)$.  That domain is a ball: for every $x$,
\[
  x_0+r\,x\in B_\rho(x_0)\iff\lvert r\,x\rvert<\rho\iff r\lvert x\rvert<\rho
  \iff\lvert x\rvert<\rho/r,
\]
using $r>0$.  Hence, writing $G_r(x):=(r/r^{q})\,Du(x_0+r\,x)$,
\begin{equation}\label{eq:rig-ur}
  u_r\ \text{is}\ p\text{-harmonic in}\ B_{\rho/r}(0)\ \text{with weak gradient}\ G_r .
\end{equation}
Moreover $G_r$ is the classical gradient of $u_r$: with $T_r(x):=x_0+r\,x$ we have
$u_r=r^{-q}\bigl(u\circ T_r-u(x_0)\bigr)$, and $u$ is differentiable at $T_r(x)$ whenever
$T_r(x)\in B_\rho(x_0)$, so clause (5) of the background proposition on the gradient under an
affine substitution gives
\begin{equation}\label{eq:rig-grad}
  Du_r(x)=r^{-q}\,r\,Du\bigl(T_r(x)\bigr)=G_r(x)\qquad\text{for }x\in B_{\rho/r}(0).
\end{equation}

\emph{Step 2: the sequence.}  Put
\[
  r_n:=\frac{\rho}{n+1}>0,\qquad v_n:=u_{r_n},\qquad g_n:=G_{r_n},\qquad
  \Omega_n:=B_{n+1}(0)\qquad(n\in\mathbb{N}).
\]
Then $\rho/r_n=n+1$, so \eqref{eq:rig-ur} reads: $v_n$ is $p$-harmonic in $\Omega_n$ with weak
gradient $g_n$.  The family $(\Omega_n)$ is increasing, since $n\le n'$ gives $n+1\le n'+1$ and
hence $B_{n+1}(0)\subseteq B_{n'+1}(0)$; and $\bigcup_n\Omega_n=\mathrm{Plane}$, because every
$z$ satisfies $\lvert z\rvert<n+1$ for $n$ large.

\emph{Step 3: the $C^0$ convergence.}  The set $B_\rho(x_0)$ is open, $x_0\in B_\rho(x_0)$, $u$
is of class $C^{q}$ on $B_\rho(x_0)$ (being of class $C^{m}$ there with $q\le m$), and $q\ge1$;
and the inductive hypothesis $T(q)$ is exactly the vanishing hypothesis
$D^{j}u(x_0)[w,\ldots,w]=0$ for $1\le j<q$ and all $w$.  So
\noderef{BlowupUniformConvergence}, applied at the order $q$, gives: for every $R>0$ and every
$\varepsilon>0$ there is $\delta>0$ with
\begin{equation}\label{eq:rig-c0}
  \lvert u_r(x)-P(x)\rvert<\varepsilon\qquad
  \text{for all }r\text{ with }0<r<\delta\text{ and all }x\text{ with }\lvert x\rvert\le R,
\end{equation}
its $P_q$ being our $P$.  Since $r_n\to0$ and $r_n>0$, for each $R>0$ and $\varepsilon>0$ there
is $N$ with $\lvert v_n(x)-P(x)\rvert<\varepsilon$ for all $n\ge N$ and all
$\lvert x\rvert\le R$; that is,
\begin{equation}\label{eq:rig-c0seq}
  v_n\to P\ \text{uniformly on}\ \overline{B_R(0)}\ \text{for every }R>0 .
\end{equation}

\emph{Step 4: eventual local uniform boundedness.}  Let $K\subseteq\bigcup_n\Omega_n=\mathrm{Plane}$
be compact.  Then $K$ is bounded, so $K\subseteq\overline{B_R(0)}$ for some $R>0$.  By the
operator-norm background proposition, $\lvert P(x)\rvert\le\lVert Q\rVert\,\lvert x\rvert^{q}$,
so $\lvert P\rvert\le B:=\lVert Q\rVert R^{q}$ on $\overline{B_R(0)}$.  By
\eqref{eq:rig-c0seq} with $\varepsilon:=1$ there is $N$ with $\lvert v_n-P\rvert<1$ on
$\overline{B_R(0)}$ for $n\ge N$, whence $\lvert v_n(z)\rvert\le B+1$ for all $n\ge N$ and
$z\in K$.  Taking $M:=B+2>0$ gives the hypothesis of eventual local uniform boundedness.

\emph{Step 5: the gradient subsequence.}  Steps 2 and 4 supply exactly the two hypotheses of
\noderef{PHarmGradientCompactness} for the increasing family $(\Omega_n)$, the functions
$(v_n)$ and the weak gradients $(g_n)$.  Its part (b) therefore gives a strictly increasing
$\varphi\colon\mathbb{N}\to\mathbb{N}$ and functions $V\colon\mathrm{Plane}\to\mathbb{R}$ and
$G\colon\mathrm{Plane}\to\mathrm{Plane}$ such that for every compact
$K\subseteq\bigcup_n\Omega_n=\mathrm{Plane}$,
\begin{equation}\label{eq:rig-sub}
  v_{\varphi(i)}\to V\quad\text{and}\quad Dv_{\varphi(i)}\to G\qquad\text{uniformly on }K .
\end{equation}
Note that no claim is made, and none is needed, relating $G$ to the gradient of $V$.

\emph{Step 6: the function limit is $P$.}  Fix $R>0$.  Since $\varphi$ is strictly increasing,
$\varphi(i)\ge i$, so $r_{\varphi(i)}\to0$; hence \eqref{eq:rig-c0} gives
$v_{\varphi(i)}\to P$ uniformly on $\overline{B_R(0)}$, exactly as in Step 3.  On the other
hand \eqref{eq:rig-sub} with $K:=\overline{B_R(0)}$ gives $v_{\varphi(i)}\to V$ uniformly on
$\overline{B_R(0)}$.  Uniform convergence implies pointwise convergence, and limits in
$\mathbb{R}$ are unique, so $V=P$ on $\overline{B_R(0)}$.  As $R>0$ was arbitrary and every
point of $\mathrm{Plane}$ lies in some $\overline{B_R(0)}$,
\begin{equation}\label{eq:rig-VP}
  V=P\quad\text{on }\mathrm{Plane}.
\end{equation}

\emph{Step 7: $P$ is $p$-harmonic on each ball.}  Fix $R>0$ and choose $I\in\mathbb{N}$ with
$I\ge R$.  For $i\in\mathbb{N}$ put
\[
  w_i:=v_{\varphi(i+I)},\qquad h_i:=g_{\varphi(i+I)} .
\]
Since $\varphi$ is strictly increasing, $\varphi(i+I)\ge i+I\ge I\ge R$, so
$\varphi(i+I)+1>R$ and therefore $B_R(0)\subseteq B_{\varphi(i+I)+1}(0)=\Omega_{\varphi(i+I)}$.
By Step 2, $w_i$ is $p$-harmonic in $\Omega_{\varphi(i+I)}$ with weak gradient $h_i$, and
$B_R(0)$ is open, so part (a) of \noderef{PHarmRestrict} gives
\[
  w_i\ \text{is}\ p\text{-harmonic in}\ B_R(0)\ \text{with weak gradient}\ h_i,
  \qquad\text{for every }i\in\mathbb{N}.
\]
Next, let $K\subseteq B_R(0)$ be compact; such a $K$ is in particular a compact subset of
$\mathrm{Plane}=\bigcup_n\Omega_n$.  The sequence $(w_i)_i$ is the tail
$i\mapsto v_{\varphi(i+I)}$ of $(v_{\varphi(i)})_i$.  The sequence $(h_i)_i$ agrees on $K$ with
the corresponding tail of $(Dv_{\varphi(i)})_i$: indeed
$K\subseteq B_R(0)\subseteq\Omega_{\varphi(i+I)}=B_{\rho/r_{\varphi(i+I)}}(0)$, which is the set
on which \eqref{eq:rig-grad} asserts $G_{r_{\varphi(i+I)}}=Du_{r_{\varphi(i+I)}}$, that is
$h_i=Dv_{\varphi(i+I)}$ on $K$.  A tail of a sequence converging uniformly on a set converges
uniformly on that set to the same limit, and two sequences agreeing on $K$ have the same uniform
limit behaviour on $K$, so by \eqref{eq:rig-sub} and \eqref{eq:rig-VP}
\[
  w_i\to P\quad\text{and}\quad h_i\to G\qquad\text{uniformly on }K .
\]
These are precisely the three hypotheses of \noderef{PHarmLimit} for the domain $B_R(0)$, the
sequences $(w_i)$, $(h_i)$ and the limits $P$, $G$.  Hence
\begin{equation}\label{eq:rig-ballharm}
  P\ \text{is}\ p\text{-harmonic in}\ B_R(0)\ \text{with weak gradient}\ G,
  \qquad\text{for every }R>0 .
\end{equation}

\emph{Step 8: $P$ is $p$-harmonic on the whole plane.}  Apply \eqref{eq:rig-ballharm} with
$R:=n+1$ for each $n\in\mathbb{N}$: $P$ is $p$-harmonic in $\Omega_n=B_{n+1}(0)$ with weak
gradient $G$, the same $G$ for every $n$.  The family $(\Omega_n)$ is increasing by Step 2, so
\noderef{PHarmUnionExhaustion} applies and gives that $P$ is $p$-harmonic in
$\bigcup_n\Omega_n$ with weak gradient $G$.  By Step 2 that union is all of $\mathrm{Plane}$, so
\begin{equation}\label{eq:rig-univ}
  P\ \text{is}\ p\text{-harmonic on all of}\ \mathrm{Plane}\ \text{with weak gradient}\ G .
\end{equation}

\emph{Step 9: Lewis rigidity.}  We have $q\ge2$, the map $Q$ is a continuous $q$-linear form on
$\mathrm{Plane}$, the function $x\mapsto Q[x,\ldots,x]$ equals $P$ by definition, and
\eqref{eq:rig-univ} says that this function is $p$-harmonic on all of $\mathrm{Plane}$ with the
weak gradient $G$.  These are exactly the hypotheses of \noderef{LewisRigidityInput}, whose
conclusion is that $Q[x,\ldots,x]=0$ for every $x\in\mathrm{Plane}$.  Note that the weak
gradient there is quantified along with $Q$, so it is enough to exhibit \emph{some} weak
gradient for $P$; we never need to know that $G$ is the gradient of $P$.

Finally, $Q[w,\ldots,w]=\tfrac1{q!}D^{q}u(x_0)[w,\ldots,w]$ by definition of $Q$ and clause (7)
of the recursion background proposition, and $q!\neq0$, so $Q[w,\ldots,w]=0$ gives
$D^{q}u(x_0)[w,\ldots,w]=0$ for every $w$.  This is \eqref{eq:rig-goal}, hence $T(q+1)$, and the
induction is complete.

\medskip\noindent Taking $q=m+1$ gives $T(m+1)$, which contains the asserted vanishing for every
$j$ with $2\le j\le m$.
\end{proof}
